\documentclass[12pt]{article}

\usepackage{amsmath, amssymb}
\usepackage{geometry}
\usepackage{setspace}
\usepackage{enumitem}

\geometry{margin=1in}
\setstretch{1.2}

\begin{document}

\begin{center}
\Large \textbf{CSE400 – Fundamentals of Probability in Computing} \\
\vspace{0.2cm}
\large \textbf{Lecture 10: Randomized Min-Cut Algorithm}
\end{center}

\vspace{0.5cm}

\section{Min-Cut Problem}

\subsection{Why Use Min-Cut?}

The \textbf{min-cut algorithm} is used in various applications to solve problems related to:
\begin{enumerate}
    \item Network connectivity
    \item Network reliability
    \item Optimization
\end{enumerate}

The motivation is explained through specific applications:

\begin{itemize}
    \item \textbf{Network Design} \\
    Min-cut helps in improving the efficiency of communication and optimizing network flow. The algorithm is used in network design to find the minimum capacity cut.
    
    \item \textbf{Communication Networks} \\
    Min-cut is useful for understanding the vulnerability of networks to failures. It helps in building robust and fault-tolerant communication networks.
    
    \item \textbf{VLSI Design} \\
    In Very Large Scale Integration (VLSI) design, min-cut is useful for partitioning circuits into smaller components, which leads to reduced interconnectivity complexity.
\end{itemize}

\section{What is Min-Cut?}

\subsection{Cut-Set}

\textbf{Definition (Cut-Set):} \\
A cut-set in a graph is a set of edges whose removal breaks the graph into two or more connected components.

\subsection{Minimum Cut (Min-Cut)}

Given:
\begin{itemize}
    \item A graph $G = (V, E)$
    \item $n = |V|$ vertices
\end{itemize}

\textbf{Definition (Min-Cut Problem):} \\
The minimum cut, or min-cut, problem is to find a minimum cardinality cut-set in the graph $G$.

That is:
\begin{itemize}
    \item Among all possible cut-sets,
    \item Choose the one with the smallest number of edges.
\end{itemize}

\subsection{Sensitivity of Min-Cut Algorithms}

Min-cut algorithms such as Karger’s algorithm are:
\begin{itemize}
    \item Randomized
    \item Sensitive to the initial choice of edges
\end{itemize}

Logical implication:
\begin{itemize}
    \item If the algorithm contracts critical edges early,
    \item Then the resulting cut may be smaller or incorrect, depending on randomness.
\end{itemize}

\section{Edge Contraction}

\subsection{Main Operation}

The main operation in the min-cut algorithm is \textbf{edge contraction}.

\textbf{Definition (Edge Contraction):} \\
Edge contraction is an operation that:
\begin{enumerate}
    \item Removes an edge from the graph, and
    \item Simultaneously merges the two vertices connected by that edge.
\end{enumerate}

\subsection{Contracting an Edge $(u, v)$}

When contracting an edge $(u, v)$:

\begin{enumerate}
    \item Vertices $u$ and $v$ are merged into a single vertex.
    \item All edges connecting $u$ and $v$ are eliminated.
    \item All other edges in the graph are retained.
\end{enumerate}

Properties of the resulting graph:
\begin{itemize}
    \item Parallel edges may exist
    \item Self-loops do not exist
\end{itemize}

\section{Successful and Unsuccessful Min-Cut Runs}

\subsection{Successful Min-Cut Run}

\textbf{Definition:} \\
A successful min-cut run refers to the success in the outcome of an algorithm designed to find the minimum cut in a graph.

Logical meaning:
\begin{itemize}
    \item The algorithm correctly identifies the true minimum cut.
\end{itemize}

\subsection{Unsuccessful Min-Cut Run}

\textbf{Definition:} \\
An unsuccessful min-cut run refers to an iteration of a min-cut algorithm where the algorithm fails to correctly identify the minimum cut of a given graph.

Logical meaning:
\begin{itemize}
    \item Due to poor random choices (e.g., contracting critical edges),
    \item The algorithm outputs a cut that is not minimum.
\end{itemize}

\section{Max-Flow Min-Cut Theorem}

\subsection{Statement of the Theorem}

\textbf{Max-Flow Min-Cut Theorem:} \\
In a flow network, the maximum amount of flow passing from the source to the sink is equal to the total weight of the edges in a minimum cut.

\subsection{Related Definitions}

Let a cut divide the vertex set into two parts $X$ and $Y$:

\begin{itemize}
    \item \textbf{Capacity of a cut} \\
    The sum of the capacities of edges oriented from a vertex in $X$ to a vertex in $Y$.
    
    \item \textbf{Minimum cut} \\
    The cut in the network that has the smallest possible capacity.
    
    \item \textbf{Minimum cut capacity} \\
    The capacity value of the minimum cut.
    
    \item \textbf{Maximum flow} \\
    The largest possible flow from the source $S$ to the sink $T$.
\end{itemize}

Logical implication of the theorem:
\begin{itemize}
    \item Computing maximum flow also determines the minimum cut capacity.
    \item Both quantities are equal.
\end{itemize}

\section{Deterministic Min-Cut Algorithm}

\subsection{Stoer–Wagner Min-Cut Algorithm}

\subsection{Theorem Used}

Let:
\begin{itemize}
    \item $s$ and $t$ be two vertices of a graph $G$
    \item $G/\{s,t\}$ be the graph obtained by merging vertices $s$ and $t$
\end{itemize}

\textbf{Statement:} \\
A minimum cut of $G$ can be obtained by taking the smaller of:
\begin{enumerate}
    \item A minimum $s$-$t$ cut of $G$, and
    \item A minimum cut of $G/\{s,t\}$
\end{enumerate}

\subsection{Reasoning Behind the Theorem}

The theorem holds because exactly one of the following cases must occur:

\begin{enumerate}
    \item There exists a minimum cut of $G$ that separates $s$ and $t$. \\
    Then the minimum $s$-$t$ cut of $G$ is a minimum cut of $G$.
    
    \item There is no minimum cut that separates $s$ and $t$. \\
    Then the minimum cut of the merged graph $G/\{s,t\}$ is a minimum cut of $G$.
\end{enumerate}

\subsection{Pseudocode}

\textbf{Algorithm 1: MinimumCutPhase$(G, a)$}

\begin{enumerate}
    \item Initialize $A \leftarrow \{a\}$
    \item While $A \neq V$:
    \begin{itemize}
        \item Add to $A$ the most tightly connected vertex
    \end{itemize}
    \item Return the cut weight as the cut of the phase
\end{enumerate}

\textbf{Algorithm 2: MinimumCut$(G)$}

\begin{enumerate}
    \item While $|V| \ge 1$:
    \begin{itemize}
        \item Choose any $a \in V$
        \item Run MinimumCutPhase$(G, a)$
    \end{itemize}
    \item If the cut of the phase is lighter than the current minimum cut:
    \begin{itemize}
        \item Store it as the current minimum cut
    \end{itemize}
    \item Shrink $G$ by merging the last two vertices added
    \item Return the minimum cut
\end{enumerate}

\section{Randomized Min-Cut Algorithm}

\subsection{Why Randomized Algorithms?}

Randomized algorithms:
\begin{enumerate}
    \item Provide a probabilistic guarantee of success
    \item Provide a more accurate estimate of the minimum cut with fewer iterations
\end{enumerate}

Additional reasons listed:
\begin{itemize}
    \item Efficiency
    \item Parallelization
    \item Approximation guarantees
    \item Avoidance of worst-case instances
    \item Heuristic nature
    \item Robustness
\end{itemize}

\section{Karger’s Randomized Algorithm}

Karger’s algorithm repeatedly:
\begin{enumerate}
    \item Randomly selects an edge
    \item Contracts the selected edge
    \item Continues until only two vertices remain
    \item The remaining edges define a cut
\end{enumerate}

The output may or may not be the true minimum cut, depending on randomness.

\subsection{Recursive Randomized Min-Cut – Pseudocode}

\textbf{Algorithm 3: RECURSIVE-RANDOMIZED-MIN-CUT$(G, \alpha)$}

\textbf{Input:}
\begin{itemize}
    \item Undirected multigraph $G$ with $n$ vertices
    \item Integer constant $\alpha > 0$
\end{itemize}

\textbf{Output:}
\begin{itemize}
    \item A cut $C$ of $G$
\end{itemize}

\textbf{Steps:}

\begin{enumerate}
    \item If $n \le 3$:
    \begin{itemize}
        \item Compute min-cut using brute-force (exhaustive) search
        \item Set $C \leftarrow$ min-cut found
    \end{itemize}
    \item Else:
    \begin{itemize}
        \item For $i = 1$ to $\alpha$:
        \begin{enumerate}
            \item Construct $G'$ by applying $n - \lceil n/\sqrt{\alpha} \rceil$ random contraction steps on $G$
            \item Compute $C' \leftarrow$ RECURSIVE-RANDOMIZED-MIN-CUT$(G', \alpha)$
            \item If $i = 1$ or $|C'| < |C|$:
            \begin{itemize}
                \item Set $C \leftarrow C'$
            \end{itemize}
        \end{enumerate}
    \end{itemize}
    \item Return $C$
\end{enumerate}

\section{Deterministic vs Randomized Min-Cut}

\subsection{Deterministic Min-Cut}

\begin{itemize}
    \item Always guarantees an exact minimum cut
    \item May have higher time complexity for large graphs
\end{itemize}

Stoer–Wagner algorithm complexity:
\[
O(V \cdot E + V^2 \log V)
\]

\subsection{Randomized Min-Cut}

\begin{itemize}
    \item Provides an approximate minimum cut with high probability
\end{itemize}

Karger’s algorithm complexity:
\[
O(V^2)
\]

\section{Theorem for Min-Cut Set}

\textbf{Theorem:} \\
The algorithm outputs a minimum cut set with probability at least:
\[
\frac{2}{n(n-1)}
\]

\section{Python Simulation}

Students are instructed to:
\begin{enumerate}
    \item Open the Campuswire post for Lecture 10
    \item Download the \texttt{.ipynb} file shared in the comments
\end{enumerate}

\vspace{0.5cm}

\begin{center}
\textbf{End of Lecture 10} \\
Thank You
\end{center}

\end{document}
